%=====================================================================
\section{What we know, and what it means}\label{sec:discussion}
\subsection{The picture}
\begin{enumerate}
  \item \textbf{Passive}: the imposed winding produces exactly $2|k|$ defects on fixed sites just inside the wall, at a depth set only by the
        imposed spacing ($\approx0.56+0.094/\rho$ cells), for every shape. Bulk charge 0.
  \item \textbf{Active circles}: two $+\half$ leave the wall and form a bulk of charge $+1$ --- the Euler characteristic of the disc,
        as if the wall had planar anchoring --- while $2k-2$ stay pinned on fixed sites. This holds for $R\approx0.75$--$2\,\la$ at every
        density the anchoring resolves, at every $k$, and is an active effect. The layer and the bulk are almost separate populations:
        wall defects stay $\gg1000$ time units, only 4--10\% of bulk $+\half$ ever dock.
  \item \textbf{The bulk is born with the vortex}: below $R/\la\approx0.4$ all $2k$ $+\half$ sit at the wall and nothing circulates; circulation,
        tilt and the $+1$ bulk appear together at $\approx0.4\,\la$ for every $k$ and $\rho$; at large $R$ the layer becomes less charged
        (bulk $>1$, lattice pinned but under-occupied) as the circulation weakens (Sec.~\ref{sec:onset}).
  \item \textbf{The circular layer is a pinned lattice tilted with the flow}: period $2k-2$, wall $+\half$ leaning 60--65$^\circ$ downstream,
        a single vortex whose circulation collapses on $R/\la$. Layer, pinning and vortex disappear together above $\approx2.5\,\la$.
  \item \textbf{Negative winding}: the wall lattice still has the period of a $+1$ bulk ($2|k|+2$) but is under-occupied, so the bulk is
        positive but below $+1$ until the wall is crowded ($|k|\gtrsim15$). Activity breaks the mirror symmetry.
  \item \textbf{Shape}: at the same size and density, elongation destroys the split. Round shapes (aspect $\lesssim1.2$) behave like
        circles; from aspect $\approx1.4$--1.5 there is no plateau, no pinned lattice and continuous wall--bulk exchange. Ovals and ellipses
        follow the same curve in the aspect ratio, convex or not, and a single lobe-shaped oval behaves like the lobes of a pinched one.
        \textbf{The split is geometric, not fixed by topology.}
  \item \textbf{Everything changes at the same aspect ratio}: charge plateau, lattice pinning, single vortex $\to$ two vortices, escape
        and docking rates (small ovals: aspect 1.17 circle-like, 1.29 intermediate, $\ge1.5$ lost).
  \item \textbf{Crowding pushes back}: a high imposed density ($k$ large at fixed perimeter) refills the wall lattice, stops the upstream
        creep and restores a plateau with a total bulk of $\approx+1$, in ovals as in the mirror runs.
\end{enumerate}
\subsection{Proposed mechanism (to be tested)}
\emph{Elastic part (Sec.~\ref{sec:theory}):} in a circle the wall layer shields the bulk exactly (2D shell theorem with same-sign images), which
explains why the bulk is blind to $k$; the same model predicts the passive layer depth without free parameters (correlation 0.996 over 310 passive
circles) and the screened approach measured in the mirror circles. In an ellipse the passive layer stays evenly spaced and the shielding is imperfect,
but the leftover field is weak ($\lesssim0.06\,\pi K$ inside), so the flow is the more likely driver of the layer loss. Negative winding confirms this (Sec.~\ref{sec:smirror}): the $-\half$ wall layer, whose defects are not self-propelled, survives elongation and the shielding weakens only modestly.\par

In a circle one vortex shears the whole wall in one sense; the tilted wall $+\half$ are held on their sites by the flow and the
anchoring, and the sites are all occupied. In an elongated cavity the flow organises into two vortices (or one that keeps reversing),
which meet along the wall; there the near-wall flow is weak and changes sign, sites empty, and wall defects creep upstream into
vacancies and escape into the bulk, while bulk $+\half$ dock. The charge is then spread through the cavity instead of being stored in a
wall layer. The vacancy--drift correlation (Spearman $-0.82$) and the coincidence of the vortex transition with the charge transition
support this; a direct test is to map vacancies and escapes along the wall relative to the vortex stagnation points.

\subsection{Status of the hypotheses}
\begin{table}[H]\centering\small
\begin{tabularx}{\textwidth}{lXX}\toprule
& Question & Status\\\midrule
H1 & Topological or geometric? & \textbf{Geometric} (Part III); ellipse eMode 0 vs 1 test not needed any more\\
H2 & Does the effect fade at low density ($\rho\la<1$)? & No, down to $\rho\la=0.44$ (Sec.~\ref{sec:ccharge})\\
H3 & Does activity enter only through $\la$? & No; $\zeta=0.07$ is the clean working point (Sec.~\ref{sec:activity})\\
H4 & Does the $+1$ bulk appear with circulation? & \textbf{Supported}: onsets coincide at $R/\la\approx0.4$ for all $k$, $\rho$, both vanish together at large $R$, and the layer is lost where the vortex splits in ovals; causality not tested\\
H5 & Negative winding & Bulk positive, $\to+1$ at $|k|\gtrsim15$; vacancies (Sec.~\ref{sec:mirror}); in ellipses the $-\half$ layer survives elongation (Sec.~\ref{sec:smirror})\\
H6 & Layer formation vs passive layer & Active layer holds $2k-2$, passive $2k$; passive depth law (Sec.~\ref{sec:passive})\\
H7 & Docking at the layer edge & Rates measured (Secs.~\ref{sec:csurv}, \ref{sec:ssurv}); pathway not analysed\\
H8 & Two-lobe cavities: topology vs circulation & Lobes do not act as separate cavities; lobe shape alone loses the layer\\\bottomrule
\end{tabularx}
\end{table}

\subsection{Open items and next runs}
\begin{itemize}
  \item \textbf{Circles to 10 seeds} (running; $+k$ $\approx$ 9--10 Oct, $-k$ $\approx$ 13--14 Oct): refresh every circle figure; check the
        ``less-charged layer'' (Sec.~\ref{sec:sweep}).
  \item \textbf{Oval sweep} (drafted, not launched): 9 shapes $\times$ 3 equivalent sizes $\times$ 4 densities, seed count to decide
        (10 per point would take $\approx4$ days on the server). Open choices: mirror subset, a high-$\rho$ column, ellipses on the server.
  \item Mechanism: vacancies and escapes along the wall vs the vortex stagnation points; near-wall flow and tilt along the arclength in
        ovals; docking pathway (H7).
  \item Fix the ellipse wall distance in \texttt{edge\_criteria.wall()} and \texttt{layer\_depth.py} before quoting $d^*$ or tilt for
        ellipses (Appendix~\ref{app:issues}).
  \item Channel geometry (the infinite-aspect limit) only if a referee asks.
\end{itemize}
\subsection{Candidate split between paper and SI}
\textbf{Main text:} passive vs active circles (maps, $Q(r)$), the layer diagram, the mirror result, the shape result ($Q(c)$, $f^*$
collapse, small-oval crossover), the mechanism figure (pinning, vacancies, drift, one vs two vortices, survival).
\textbf{SI:} resolution limit, $\la$ calibration, compressibility, every edge criterion and their comparison, pair correlations,
activity scan, Fourier details, the $k$ scan, pinch-off and lobe reference, the label tracking. Marked ``SI?'' in the margin above.
